%==============================================================================
% MATH 5080
% Class:        13
% Date:         Oct 13th
%==============================================================================

\subsection{L'Hopital's Rule}

\begin{theorem}[L'Hopital's Rule, 0/0 case]
    \label{thm:lhopital-1}
    Let $f, g$ be differentiable on 
    $U = N^*(c, \delta)$ 
    for some $\delta > 0$, s.t. 
    \begin{enumerate}
        \item $\lim_{x \to c}f(x) = \lim_{x \to c}g(x) = 0$
        \item $g^{\prime}(x) \neq 0, \forall x \in U$
        \item $\lim_{x \to c} \frac{f^{\prime}(x)}{g^{\prime}(x)} = L \in \mathbb{R} \cup \{ \pm \infty \} $
    \end{enumerate}
    Then
    \[
        \lim_{x \to c} \frac{f(x)}{g(x)} = L
    \]
\end{theorem}

\begin{remark}
    The above conditions are sufficient for $g(x) \neq 0, \forall x \in U$.
    ATC otherwise, we can extend $g: N^*(x, \delta) \to \mathbb{R}$
    to $\tilde{g}: N(x, \delta) \to \mathbb{R}$ by letting 
    \[
        \tilde{g} (c) = \lim_{x \to c}g(x) = 0
    \]
    If $g(x) = 0$, then MVT on $[x, c]$ 
    implies $\exists c_1 \in (x, \c)$ s.t.
    \[
        g(c_1)= \frac{
            \tilde{g}(c) - \tilde{g}(x)
        }{
            c - x
        } = 0
    \]
    which is a contradiction.
\end{remark}

\begin{lemma}[Cauchy Mean Value Theorem]
    Let $f, g $ be continuous on $[a,b]$, 
    differentiable on $(a,b)$, then 
    $\exists c\in (a,b)$ s.t.
    \[
        \left( f(b) - f(a) \right) \cdot g^{\prime}(c)
        =
        \left( g(b) - g(a) \right) \cdot f^{\prime}(c)
    \]
\end{lemma}

\begin{remark}
    \hfill
    \begin{enumerate}
        \item $g(x) = x$ recovers MVT.
        \item This is not a direct consequence of $MVT$
    \end{enumerate} 
\end{remark}

\begin{proof}
    Define $h(x)$
    \[
        h(x) = 
        \left( f(b) - f(a) \right) g(x) - 
        \left( g(b) - g(a) \right) f(x)
    \]
    then $h$ is continuous on $[a, b]$, 
    differentiable on $(a,b)$,
    and $h(a) = f(b)g(a) - g(b)f(a) =h(b)$.
    By Rolle's $\exists c \in (a, b)$ s.t. $h^{\prime}(c) = 0$,
    i.e.
    \[
        \left( f(b) - f(a) \right) \cdot g^{\prime}(c)
        =
        \left( g(b) - g(a) \right) \cdot f^{\prime}(c)
    \]
\end{proof}

\begin{proof}[Proof of 
    \ref{thm:lhopital-1}
]
\begin{itemize}
    \item Case $L \in \mathbb{R}$: Suffices to show 
    \[
        \lim_{x \to c^{-}} \frac{f(x)}{g(x)} = L
    \]
    The idea is to
    \begin{enumerate}
        \item Use Cauchy MVT to write 
        \[
            \frac{{f}(x)}{{g}(x)}
            = \frac{f^{\prime}(c_x)}{g^{\prime}(c_x)}
        \]
        \item Taking the limit on both sides and show that 
        the limit is equal to $L$.
    \end{enumerate} 

    \textbf{Step 1}: Fix any $x < c, x \in U$.
    Since we know that $f, g$ differentiable and therefore
    continuous on $[x, c)$, and 
    $\lim_{x \to c} f(x) = \lim_{x \to c}g(x) = 0$,
    we can define 
    $\tilde{f}(c) = \tilde{g}(c) = 0$ and extend
    $f, g$ to $\tilde{f}, \tilde{g}$ on $[x, c]$.
    We have $\tilde{f}, \tilde{g}$ differentiable on $(x, c)$. 
    \medskip

    Applying Cauchy MVT, $\exists c_x \in (x, c)$ s.t. 
    \[
        \left( \tilde{f}(c) - \tilde{f}(x) \right) g^{\prime}(c_x) = 
        \left( \tilde{g}(c) - \tilde{g}(x) \right) f^{\prime}(c_x)
    \]
    By assumption, $g^{\prime} \neq 0$, which implies
    \begin{itemize}
        \item $g^{\prime}(c_x) \neq 0$
        \item $\tilde{g}(c) - \tilde{g}(x) \neq 0$. Otherwise,
        by MVT, $\exists c_2 \in (x, c)$ s.t. 
        \[
            g^{\prime}(c_2) = \frac{\tilde{g}(c)- \tilde{g}(x)}{c - x}= 0
        \]
    \end{itemize} 
    Hence,
    \[
        \left( \tilde{f}(c) - \tilde{f}(x) \right) g^{\prime}(c_x) = 
        \left( \tilde{g}(c) - \tilde{g}(x) \right) f^{\prime}(c_x)
        \implies
        \frac{\tilde{f}(c) - \tilde{f}(x)}{\tilde{g}(c) - \tilde{g}(x)}
        = \frac{f^{\prime}(c_x)}{g^{\prime}(c_x)}
    \]
    Since $\tilde{f}(c) = \tilde{g}(c) = 0$, we have
    \[
        \frac{f(x)}{g(x)}
        = \frac{f^{\prime}(c_x)}{g^{\prime}(c_x)}
    \]
    \medskip

    \textbf{Step 2}: by assumption, we have 
    \[
        \lim_{x \to c^{-}} \frac{f^{\prime}(x)}{g^{\prime}(x)} = L
    \]
    hence $\forall \ep, \exists \delta$ s.t. 
    \[
        x \in (c - \delta, c) \implies 
        \left| 
            \frac{f^{\prime}(x)}{g^{\prime}(x)} - L 
        \right|  < \ep
    \]
    For any such $x$, since $x < c_x < c$, therefore
    $ c- \delta < c_x < c$. Hence
    \[
        \left| 
            \frac{f^{\prime}(c_x)}{g^{\prime}(c_x)} - L
        \right|  < \ep
    \]
    Therefore 
    \[
        \lim_{x \to c} \frac{f^{\prime}(c_x)}{g^{\prime}(c_x)} = L 
        \implies
        \lim_{x \to c} \frac{f(x)}{g(x)} = L 
    \]
    
    \item Case $L = \pm \infty$,
    WTS 
    \[
        \lim_{x \to c^{+}} \frac{f(x)}{g(x)} = \infty
    \]
    As before, we get
    \[
        \frac{f(x)}{g(x)} = \frac{f^{\prime}(c_x)}{g^{\prime}(c_x)}
        \text{ for some } c_x \in (c, x)
    \]
    By assumption, 
    \[
        \lim_{x \to c} \frac{f^{\prime}(x)}{g^{\prime}(x)} = \infty
    \]
    implies $\forall M > 0, \exists \delta > 0$ s.t. 
    \[
        x \in (c, c + \delta) \implies \frac{f^{\prime}(x)}{g^{\prime}(x)} > M
    \]
    Hence $\forall x \in (c, c + \delta)$, 
    $c < c_x < c+ \delta$, hence 
    \[
        \frac{f(x)}{g(x)} = \frac{f^{\prime}(c_x)}{g^{\prime}(c_x)}
        > M
    \]
    which implieas 
    \[
        \lim_{x \to c^{+}} \frac{f(x)}{g(x)} = \infty
    \]
\end{itemize} 
\end{proof}

\begin{example}
    \begin{enumerate}
        \item $\lim_{x \to 0} \frac{e^{x} - 1}{x}$
        Check the conditions
        \begin{itemize}
            \item $\lim_{x \to 0} e^{x} - 1 = 0$, $\lim_{x \to 0}x  = 0$
            \item $x^{\prime} = 1 \neq 0$
            \item $f^{\prime} = e^{x}, g^{\prime} = 1$, 
            $\lim_{x \to 0} \frac{f^{\prime}}{g^{\prime}} = \frac{e^{x}}{1} = 1$
        \end{itemize} 
        Hence, L'Hopital gives 
        \[
            \lim_{x \to 0} \frac{e^{x} - 1}{x} = 1
        \]
        \item $\lim_{x \to 0} \frac{\cos x}{x}$
        Checking the conditions, $\lim_{x \to 0}\cos x = 1 \neq 0$,
        condition 1 fails.
    \end{enumerate} 
\end{example}

\begin{theorem}[L'Hopital's Rule for $x \to \infty$]
    \label{thm:lhopital-2}
    Let $f, g$ be differentiable on $(a, \infty)$ s.t. 
    \begin{enumerate}
        \item $\lim_{x \to \infty}f(x) = \lim_{x \to \infty} g(x) = 0$
        \item $g^{\prime}(x) \neq 0, \forall x \in (a, \infty)$
        \item $\lim_{x \to \infty} \frac{f^{\prime}(x)}{g^{\prime}(x)} = L \in \mathbb{R} \cup \{ \pm \infty \} $
    \end{enumerate} 
    Then,
    \[
        \lim_{x \to \infty} \frac{f(x)}{g(x)} = L
    \]
\end{theorem}

\begin{remark}
    Again $g(x) \neq 0$ is implicitly given. 
    \begin{itemize}
        \item if a single point $g(x_0) = 0$,
        and $g(x) \neq 0, \forall x \neq x_0$,
        then we can consider the theorem on
        $(x_0, \infty)$.
        \item if a sequence $\{ a_n \} $ s.t. 
        $f(a_n) = 0$, then we can find a contradiction
        with MVT 
        \remind{do we need sequence to find contradiction??
        or is it just two points that's needed??
        }
    \end{itemize} 
\end{remark}

\begin{example}
    \remind{add in example}
\end{example}

\begin{proof}
    \textit{The idea is to use 
    change of variables
    and the first form of L'Hopital.
    }

    Let $x = \frac{1}{t}$, then 
    $x \to \infty \iff t \to 0^{+}$.
    Define 
    $F(t) = f \left( \frac{1}{t} \right) $,
    $G(t) = g \left( \frac{1}{t} \right) $.

    We first show that satisfying the 
    conditions of \ref{thm:lhopital-2}
    implies the conditions of 
    \ref{thm:lhopital-1}.
    \begin{enumerate}
        \item By assumption, we have 
        $\lim_{x \to \infty}f(x) = \lim_{g \to \infty}g(x) = 0$
    \begin{align*}
        &\lim_{x \to \infty}f(x) = \lim_{g \to \infty}g(x) = 0 \\
        \implies &
        \lim_{t \to 0^{+}}f \left( \frac{1}{t} \right)  = 
        \lim_{g \to 0^{+}}g \left( \frac{1}{t} \right)  = 0 \\
        \implies &
        \lim_{t \to 0^{+}}F \left( {t} \right)  = 
        \lim_{g \to 0^{+}}G \left( {t} \right)  = 0 \\
    \end{align*}
        \item By assumption, $g^{\prime}(x) \neq 0$.
        Consider $G^{\prime}(t)$
        \begin{align*}
            &G^{\prime}(t) = g^{\prime} \left( \frac{1}{t} \right) 
            \cdot \frac{-1}{t^{2}}  \\
            \implies& G^{\prime}(t) \neq 0
        \end{align*}
        \item By assumption, 
        $\lim_{x \to \infty} \frac{f^{\prime}(x)}{g^{\prime}(x)} = L$.
        Then, by change of variables, we have
        \[
            \lim_{t \to 0^{+}} 
            \frac{f^{\prime} \left( \frac{1}{t} \right) }{
                g^{\prime} \left( \frac{1}{t} \right) 
            }
            = \lim_{x \to \infty}
            \frac{f^{\prime}(x)}{g^{\prime}{x}} = L
        \]
        By Chain Rule, 
        \begin{align*}
            F^{\prime}(t) = f^{\prime} \left( \frac{1}{t} \right) 
            \cdot \left( - \frac{1}{t^{2}} \right)  \\
            G^{\prime}(t) = g^{\prime} \left( \frac{1}{t} \right) 
            \cdot \left( - \frac{1}{t^{2}} \right)  \\
        \end{align*}
        Hence
        \begin{align*}
            \lim_{x \to 0^{+}} 
            \frac{F^{\prime}(t)}{G^{\prime}(t)} 
            =
            \lim_{x \to 0^{+}} 
            \frac{
                f^{\prime} \left( \frac{1}{t} \right) 
            }{
                g^{\prime} \left( \frac{1}{t} \right) 
            }  = L
        \end{align*}
    \end{enumerate} 
    Having verified all 3 conditions, 
    apply \ref{thm:lhopital-1}
    to $F, G$, and we get 
    \[
        \lim_{t \to 0^+}
        \frac{F(t)}{G(t)} 
        = 
        \lim_{t \to 0^+}
        \frac{F^{\prime}(t)}{G^{\prime}(t)} 
        = L
    \]
    Applying change of variables again, we get
    \[
        \lim_{x \to [infty]}
        \frac{f(x)}{g(x)} 
        =
        \lim_{t \to 0^+}
        \frac{F(t)}{G(t)} 
        = L
    \]
    
\end{proof}

\begin{remark}
    The first two L'Hopital's Rule theorems 
    (
        \ref{thm:lhopital-1},
        \ref{thm:lhopital-2}
    )
    theoretically
    cover all possible cases of limits of indeterminate
    forms. However, for convenience, we extend L'Hopital's to 
    $ \frac{\infty}{\infty}$ forms directly.
\end{remark}

\begin{theorem}[L'Hopital's Rule for $ \frac{\cdot}{\infty}$]
    \label{thm:lhopital-3}
    Let $f, g$ be differentiable on 
    $U = N^{*}(c, \delta)$ for some $\delta > 0$
    s.t. 
    \begin{enumerate}
        \item $\lim_{x \to c} g(x) = \pm \infty$
        \item $g^{\prime}(x) \neq 0, \forall x \in U$
        \item $\lim_{x \to c} \frac{f^{\prime}(x)}{g^{\prime}(x)} = L$
        $\in \mathbb{R} \cup \{ \pm \infty \} $
    \end{enumerate} 
    Then 
    \[
        \lim_{x \to c} \frac{f(x)}{g(x)} = L
    \]
\end{theorem}

\begin{remark}
    Note that \ref{thm:lhopital-3}
    does not require $\lim_{x \to c}f(x) = \pm \infty$.
    This makes the theorem stronger, 
    but means we cannot prove using 
    change of variables and \ref{thm:lhopital-1}.
\end{remark}


\begin{theorem}
    [L'Hopital's Rule for $ \frac{\cdot}{\infty}, x \to \infty$]
    \label{thm:lhopital-4}
    Let $f, g$ be differentiable on $(a, \infty)$
    s.t. 
    \begin{enumerate}
        \item $\lim_{x \to \infty}g(x) = \pm \infty$
        \item $g^{\prime}(x) \neq 0, \forall x \in (a, \infty)$
        \item $\lim_{x \to \infty} \frac{f^{\prime}(x)}{g^{\prime}(x)} = L$
        $\in \mathbb{R} \cup \{ \pm \infty \} $
    \end{enumerate} 
    then,
    \[
        \lim_{x \to \infty} \frac{f(x)}{g(x)} = L
    \]
\end{theorem}

\begin{remark}
    \ref{thm:lhopital-4} follows from
    \ref{thm:lhopital-3} by change of variables.
\end{remark}

\subsection{Taylor's Theorem}

\begin{theorem}[Taylor's Theorem]
    \label{thm:taylor}
    Let $f^{(i)}$ be continuous on $[a,b]$,
    differentiable on $(a,b)$,
    $\forall 0 \leq i \leq n$.
    Let $x_0 \in [a, b]$. Then
    $\forall x \in [a, b]$ with
    $x \neq x_{0}$, 
    $\exists c $ between $x$ and $x_{0}$ s.t. 
    \begin{align*}
        f(x) &= f(x_0) 
        + f^{\prime}(x_0) (x-x_{0})
        + \frac{
            f^{(2)}(x_0)
        }{
            2!
        }
        (x - x_{0})^{2} 
        + \hdots 
        + \frac{
            f^{(n)}(x_0)
        }{
            n!
        }
        (x - x_{0})^{n}  \\
        &
        + \frac{
            f^{(n + 1)}(c)
        }{
            n!
        }
        (x - x_{0})^{n} 
    \end{align*}
\end{theorem}

\remind{
    why is it always
    continuous on closed interval and
    differentiable on open interval??
}

\begin{remarks}
    \begin{itemize}
        \item 
    Since the theorem requires 
    $f^{(n)}$ differentiable, 
    the RHS ($f^{(n + 1)}$) is well defined.
        \item Since $c$ depends on $x$,
    the RHS is not in a polynomial form.
        \item If $n = 0$, we recover MVT.
    \end{itemize} 
\end{remarks}

\begin{proof}
    Fix $x \in [a,b], x \neq x_{0}$. 
    Let $M$ be the unique solution of 
    \begin{align*}
        f(x) &= f(x_0) 
        + f^{\prime}(x_0) (x-x_{0})
        + \frac{
            f^{(2)}(x_0)
        }{
            2!
        }
        (x - x_{0})^{2} 
        + \hdots 
        + \frac{
            f^{(n)}(x_0)
        }{
            n!
        }
        (x - x_{0})^{n} 
        + 
        M
        (x - x_{0})^{n} 
    \end{align*}.
    WTS $M = \frac{f^{(n+1)}(c)}{(n+1)!}$ for some $c$
    between $x, x_0$. 
    Define 
    \[
        F(t) = f(t) + f^{\prime}(t) (x-t)
        + \hdots + \frac{f^{(n)}(t)}{n!} (x - t)^{n}
        + M (x-t)^{n+1}
    \]
    Then $F$ is continuous on $[a,b]$,
    differentiable on $(a,b)$.
    Observe that 
    \[
        F(x) = f(x), F(x_{0}) = f(x)
    \]
    By Rolle's theorem,
    $\exists c$ between $x, x_0$ s.t. 
    $F^{\prime}(c) = 0$.
    Differentiating $F$, we get
    \begin{align*}
        F^{\prime}(c) 
        =& f^{\prime}(c) 
        + \left[ 
            - f^{\prime}(c)
            + f^{(2)}(c) (x - c)
        \right] \\
        +& \left[ 
            - 2 \frac{f^{(x)}(c)}{2!} (x-c) 
            + \frac{f^{(3)}(c)}{2!} (x-c)^2
        \right]  \\
        +& \left[ 
            - 3 \frac{f^{(3)}(c)}{3!} (x-c)^2
            + \frac{f^{(4)}(c)}{3!} (x-c)^3
        \right]   \\
        +& \hdots \\
        +& \left[ 
            - n \frac{f^{(n)}(c)}{n!} (x-c)^n
            + \frac{f^{(n+1)}(c)}{n!} (x-c)^n
        \right] \\
        +& (-M) (n+1) (x-t)^n
    \end{align*}
    Setting this to 0, we get
    \[
        M = \frac{f^{(n+1)}(c)}{(n+1)!}
    \]
\end{proof}

\begin{definition}[Taylor's Polynomials]
    Let $f^{(i)}$ be continuous on $[a,b]$,
    differentiable on $(a,b)$, $\forall 0 \leq i \leq n$>
    Let $x_0 \in [a, b]$.
    Fix any $x \in [a,b]$, $x \neq x_0$,
    then the \textbf{Taylor's Polynomials}
    for $f$ at $x_0$ are
    \begin{align*}
        p_0(x) &= f(x_0) \\
        p_1(x) &= f(x_0) + f^{\prime}(x_0)(x-x_0)\\
        p_2(x) &= f(x_0) + f^{\prime}(x_0)(x-x_0)
        + \frac{f^{(2)}(x)}{2!} (x-x_0)^2 
        \\
        &\vdots
        p_2(x) &= f(x_0) + f^{\prime}(x_0)(x-x_0)
        + \frac{f^{(2)}(x)}{2!} (x-x_0)^2 
        + \hdots
        _ \frac{f^{(n)}(x)}{n!} (x-x_0)^n
    \end{align*}
\end{definition}


\begin{remark}
    What is the relationship between $f$ and 
    $\{ p_n \} $
    \begin{itemize}
        \item If $f$ is a polynomial, then $\exists N$
        s.t. $p_n = f, \forall n \geq N$
        \item For general $f$, $p_n(x)$ is an approximation
        \item of $f$ with an upper bound on the magnitude
        of error
        \[
            \frac{
                \max \left| f^{(n+1)} \right| 
            }{
                (n+1)!
            } 
            \left| x - x_0 \right|^{n + 1}
        \]
        
    \end{itemize} 
\end{remark}

